% HD-N14. 21 November 2026. GPT-6 Astra Pro, inherited effort.
\section{Exact trace ranks for Weil eightfold characters and two CM product recipes}

\paragraph{HD-N14 (the character ranks and their functorial consequence).}
Let $A$ be an abelian eightfold of Weil type for an imaginary
quadratic field, with balanced signature $(4,4)$ and compatible
polarization $\theta$. Let $t$ be a nonzero rational Weil class,
let $\lambda\ne0$, and let
\[
 v=e^B\bigl(P(\theta)+\lambda t\bigr),
 \qquad P\in\mathbb Q[z],\quad \deg P\le8,
 \quad B\in\operatorname{NS}(A)_{\mathbb Q}.
                                                               \tag{N14.1}
\]
For contraction $c_v^j:HH^j(A)\to HH_{-j}(A)$ with this whole
character, one has
\[
 \operatorname{rank}c_v^1=16,\qquad
 \operatorname{rank}c_v^2=56+16\rho\in\{56,72,88,104\},      \tag{N14.2}
\]
where $\rho$ is the rank of the three vectors
$P,P',P''$ after truncation modulo $z^7$.
Consequently neither of the following actual CM product
objects, nor any Fourier--Mukai image of it, has whole character
of the form(N14.1):
\begin{enumerate}
\item the external product of ideal sheaves of a point on
four abelian surfaces;
\item the external product of two genus-three quotient
Abel-curve ideals of HD-N02 and the ideal sheaf of a point
on an abelian surface.
\end{enumerate}
The same conclusions hold after finite etale isogeny pullback
with its actual deck-group descent structure and conjugation
of that structure through the Fourier--Mukai equivalence.
Every factor may be chosen of CM type, so this tests concrete
replacement-object recipes at the CM points relevant here.

\paragraph{Balanced coordinates and an exact weight decomposition.}
Choose holomorphic coordinates $a_1,\ldots,a_4$ in the
$\sigma$-eigenspace and $b_1,\ldots,b_4$ in the
$\bar\sigma$-eigenspace. The compatible nondegenerate
polarization identifies dual antiholomorphic bases
$y_1,\ldots,y_4$ and $z_1,\ldots,z_4$ of the opposite
$K$-types, with
\[
 \theta=\sum_i(a_i\wedge y_i+b_i\wedge z_i),\qquad
 \tau_+=a_1a_2a_3a_4z_1z_2z_3z_4,
 \quad\tau_-=b_1b_2b_3b_4y_1y_2y_3y_4.                   \tag{N14.3}
\]
Wedge symbols are suppressed in the volume monomials.
A nonzero rational Weil class has two nonzero conjugate
coefficients, so write $\lambda t=u\tau_++u'\tau_-$ with
$uu'\ne0$. Assign weights $+1$ to $a_i,z_i$ and $-1$ to
$b_i,y_i$. This is the explicit action of the complexified
norm-one torus: $\theta$ has weight zero and $\tau_\pm$
have weights $\pm8$.

The sixteen degree-one Hochschild operators split as
\[
 U_- =\langle\iota_{a_i},y_i\wedge\rangle_{i=1}^4,
 \qquad
 U_+ =\langle\iota_{b_i},z_i\wedge\rangle_{i=1}^4.
                                                               \tag{N14.4}
\]
Each space has dimension eight. The first acts only on
$\tau_+$ and acts injectively there: its four contractions
and four wedges yield eight different monomials. The second
has the analogous property for $\tau_-$. Their Weil outputs
have weights $+7$ and $-7$, whereas their polynomial outputs
have weights $-1$ and $+1$. Thus no polynomial term cancels
the two injective Weil maps, proving the rank16 assertion.

In degree two,
\[
 HH^2(A)=\bigwedge^2U_-\ \oplus\ (U_-\otimes U_+)\
                                      \oplus\bigwedge^2U_+.
                                                               \tag{N14.5}
\]
The two outside summands have dimension28 each. On $\tau_+$,
the first has six double contractions, sixteen contraction/wedge
operators, and six double wedges. All twenty-eight resulting
monomials are distinct; the three groups have respectively
two, three and four holomorphic $a$ factors. Hence this map
is injective. Its output weight is $+6$, while the polynomial
part of the same summand has weight $-2$.
The last summand is injectively detected by its weight $-6$
component on $\tau_-$, with polynomial output of weight $+2$.
These four weights are distinct. They are also distinct from
the weight-zero image of the middle summand.
It follows that the outside summands contribute exactly56
to the rank, independently of every coefficient of $P$.

\paragraph{The middle block and its Hankel matrix.}
The mixed summand $U_-\otimes U_+$ annihilates both
$\tau_+$ and $\tau_-$. For every pair $(i,j)$, its four
operators act on $P(\theta)$ as
\[
 y_i\wedge z_j\wedge P(\theta),\qquad
 \pm y_i\wedge z_j\wedge P'(\theta),\qquad
 \pm y_i\wedge z_j\wedge P'(\theta),\qquad
 \pm y_i\wedge z_j\wedge P''(\theta).                    \tag{N14.6}
\]
The signs depend only on operator ordering. Thus there are
three independent polynomial columns to consider for each
of the sixteen pairs. The powers
$y_i z_j\theta^k$, $0\le k\le6$, provide an injective copy
of the same sixteen-dimensional space in each indicated
bidegree. Explicitly, given a coefficient of
$\eta\in\langle y_i z_j\rangle$, choose $k$ other paired
indices disjoint from those two. The coefficient of the
monomial with exactly those $k$ holomorphic factors in
$\eta\theta^k$ recovers the chosen coefficient of $\eta$,
up to a nonzero factorial. Such a choice is possible because
$k\le6$. This proves injectivity without a dimension estimate.
Different powers have different bidegrees.

Writing $P(z)=\sum_{j=0}^8p_jz^j$ and $m_j=j!p_j$, multiplying
the $k$th row by $k!$ identifies each three-column map with
\[
 M_P=\begin{pmatrix}
 m_0&m_1&m_2\\
 m_1&m_2&m_3\\
 \vdots&\vdots&\vdots\\
 m_6&m_7&m_8
 \end{pmatrix}.                                          \tag{N14.7}
\]
The full mixed block is sixteen copies of this map. Its rank
is therefore $16\rho$, where $\rho=\operatorname{rank}M_P$.
Combining with the distinct-weight outside blocks proves the
second formula in(N14.2).
Finally, multiplication by $e^B$ conjugates every degree-one
operator by the invertible shear
$(v,w)\mapsto(v,w+\iota_vB)$, as in(N11.4)--(N11.5).
Consequently neither rank changes under the full B-field
twist in(N14.1). This proves both identities.\hfill$\square$

\paragraph{All trace degrees for the point ideal on a surface.}
Let $S$ be an abelian surface and $I_p$ the ideal of a point.
It is simple. Riemann--Roch gives
$\chi(I_p,I_p)=-2$, and surface Serre duality gives
\[
 (\operatorname{ext}^0,\operatorname{ext}^1,
                   \operatorname{ext}^2)(I_p,I_p)=(1,4,1).
                                                               \tag{N14.8}
\]
There are no negative or higher self-Ext groups.
The degree-zero trace sends the identity to
$1-[p]$, so is injective. The degree-two trace is an
isomorphism to $H^2(\mathcal O_S)$, by Serre duality with
the nonzero identity evaluation in degree zero.
In degree one, multiplication of $1-[p]$ by
$H^1(\mathcal O_S)$ detects its two-dimensional $H^{0,1}$
block, and contraction of $-[p]$ by $H^0(T_S)$ detects
its two-dimensional $H^{1,2}$ block. Hence $c_{I_p}^1$
has rank four, and the four-dimensional source self-Ext
space in(N14.8) has injective trace as well.
Thus every degree of the trace is injective for this object.

For the ideal $I_p$ of a point on an abelian threefold, the
exact sequence $I_p\to\mathcal O\to\mathcal O_p$ and the
local Koszul resolution give
\[
 0\longrightarrow T_pX\longrightarrow
 \operatorname{Ext}^1(I_p,I_p)\longrightarrow
 H^1(\mathcal O_X)\longrightarrow0.
\]
Here $\operatorname{Hom}(I_p,\mathcal O_p)=T_pX$,
$\operatorname{Ext}^1(I_p,\mathcal O_X)=H^1(\mathcal O_X)$,
and the latter group's map to
$\operatorname{Ext}^1(I_p,\mathcal O_p)$ is zero, since it is
the restriction of the scalar coherent-cohomology class to
the point. These facts give the displayed sequence directly
from the long exact self-Ext sequence. Thus its degree-one
self-Ext dimension is six, and Serre duality gives the full
array $(1,6,6,1)$. Its character $1-[p]$ has degree-one
contraction rank six: the three scalar cohomology directions
and the three contractions of the point volume occupy distinct
Hodge summands. Therefore evaluation and trace in degree one
are injective between six-dimensional spaces. The degree-two
trace is the transpose of degree-one evaluation in dimension
three, so it too is injective. The identity/Serre-dual identity
arguments treat degrees zero and three. Consequently the
threefold point ideal has injective trace in every degree.

The genus-three quotient Abel-curve ideals $J$ of HD-N02
have
\[
 (\operatorname{ext}^0,\operatorname{ext}^1,
        \operatorname{ext}^2,\operatorname{ext}^3)(J,J)
                    =(1,6,6,1).                           \tag{N14.9}
\]
The degree-one and degree-two trace injections were proved
there. In degree zero the rank-one identity character is
nonzero, and in degree three the Serre-dual identity
calculation proves injectivity. Hence these objects too
have injective trace in every degree.

\paragraph{The actual product and Fourier trace maps.}
For external products of these objects, the derived Kunneth
isomorphism identifies every self-Ext group with the direct
sum over multi-degrees of tensor products of the factor
self-Ext groups. The trace on each summand is the tensor
product of the factor traces, with the usual Koszul sign.
The Hochschild Kunneth decomposition separates these
multi-degrees, so every product trace is injective.
For four point ideals the dimensions in degrees one and two
are therefore
\[
 \operatorname{ext}^1=4\cdot4=16,\qquad
 \operatorname{ext}^2=4+\binom42\,4^2=100.                 \tag{N14.10}
\]
For two genus-three ideals and one surface point ideal they are
\[
 \operatorname{ext}^1=6+6+4=16,\qquad
 \operatorname{ext}^2=6+6+1+6\cdot6+6\cdot4+6\cdot4=97.
                                                               \tag{N14.11}
\]

There is a dimension-eight duality step here that must be
stated explicitly. On an eightfold, $\sigma^6$ is the transpose
of $\operatorname{ev}^2$, under the perfect pairings between
$\operatorname{Ext}^2$ and $\operatorname{Ext}^6$ and between
the corresponding Hochschild groups. It is \emph{not} the
self-transpose degree-two assertion used on the fourfold in
HD-N11. The same trace/Atiyah contraction identity(N11.3),
now with total Serre degree eight, proves this statement.
Since $\sigma^6$ is injective for the products above,
$\operatorname{ev}^2$ is surjective. Since $\sigma^2$ is also
injective, their composite consequently has rank
\[
 \operatorname{rank}c_F^2=\operatorname{ext}^2(F,F)
                         \in\{100,97\}.                   \tag{N14.12}
\]
An actual Fourier--Mukai equivalence gives isomorphisms of
self-Ext and Hochschild groups and the commutative trace
and evaluation diagrams. Hence it preserves the ranks in
(N14.12), including both trace injections used to prove them.
Neither100 nor97 belongs to the list(N14.2). This proves the
claimed obstruction for these product/equivalence recipes.

For an isogeny $q:B\to A$, the pullback $q^*F$ has its actual
$\ker(q)$-linearization and descends to $F$.
The invariant self-Ext groups are identified with those of
$F$ by finite etale descent; the pullback map on Hochschild
homology is an isomorphism over $\mathbb C$.
Thus the same injections hold on every invariant self-Ext
space. Restricting the degree-eight Serre pairing to invariants
again shows that ambient $\operatorname{ev}^2$ surjects onto
invariant degree-two self-Ext, and its trace composite has
rank100 or97. Conjugating these actual functors and their
coherent action by a Fourier--Mukai equivalence preserves
this statement. No unconstructed comparison of equivariant
categories is needed.

\paragraph{Exhaustion of the explicit basic-seed list.}
A line bundle on an abelian variety of dimension $g$ has
self-Ext algebra $H^*(\mathcal O)$, degree-one dimension $g$,
and injective trace in all degrees: multiplication by its
nonzero exponential is invertible on the scalar cohomology
summands. A reduced point sheaf has self-Ext algebra
$\bigwedge^*T_pX$, again degree-one dimension $g$.
Contraction of its top point class injects the degree-$j$
Koszul self-Ext space into $H^{g-j,g}(X)$; thus all its
traces are injective as well.

Consider any product of line bundles or reduced point sheaves
on positive-dimensional abelian factors, point ideals on
surfaces or threefolds, and the proved genus-three Abel ideals,
whose factor dimensions sum to eight. Every trace is injective
by the same Kunneth maps. In dimension eight the degree-seven
trace is the transpose of degree-one evaluation, so its
injectivity makes $\operatorname{rank}c^1$ equal to
the sum of the degree-one self-Ext dimensions of the factors.
A line or point sheaf contributes $g$, while each listed
ideal contributes $2g$. Formula(N14.2) requires rank16,
so no positive-dimensional line-bundle or point-sheaf factor
can occur. The remaining dimensions satisfy
\[
 2a+3b=8,\qquad a,b\ge0.
\]
The only solutions are $(a,b)=(4,0)$ and $(1,2)$.
Their self-Ext dimensions are exactly100 and97 as above,
because threefold point ideals have the same array(N14.9).
Thus the obstruction exhausts this explicit basic-seed list,
including shifts, factor line-bundle twists and the actual
isogeny/Fourier transports already described.

\paragraph{Certificate and scope.}
The exact integer exterior-algebra certificate
\nolinkurl{eightfold_character_certificate.py} checks all16
degree-one Weil columns, all56 outside degree-two columns,
and all64 mixed columns coefficient by coefficient for the
nine independent moments of $P$. Its112 coefficient-recovery
rows verify the stated middle-block injection in every degree.
The twenty-one additional full-map rank calculations modulo1009
are supplementary checks of the universal identities.
The categorical trace and transport statements have the
separate proofs above. These statements concern the displayed
actual products and their isogeny/Fourier transports. They do
not classify cones, extensions coupling the factors, or other
perfect complexes on a CM eightfold.
